% !TeX root = ../../Book2.tex
% 纲要标题：### 18.4 交叉积与相对 Brauer 群
% 纲要位置：计划纲要.md，第 1652 行。

\section{交叉积与相对 Brauer 群}
\label{b2:sec:1804}

循环代数只需一个生成自同构和一个参数。一般有限 Galois 扩张有多个自同构；记录它们对应元素相乘时的标量误差，就得到\(2\)-余圈。本节用显式方程定义这些数据，证明它们与相对 Brauer 群的对应；第四卷再从一般上链复形研究群上同调。

% 纲要要点（供目录核对）：
%
% * 交叉积与2-余圈；结合律产生余圈方程，重选基元产生余边；正式证明因子集不依赖标量提升及矩阵模型
% * 有限 Galois 扩张的相对 Brauer 群及循环情形的范数商；对角幂等元对齐两份子域作用，取角证明群运算，同类同次数及 Skolem–Noether 证明单射，明确下降因子集的逆类约定
% * 与第四卷群上同调的衔接；循环余圈化为进位形式，正文证明交叉积约化迹只由单位指标分量贡献
% * Brauer 群不能与第五卷 Brauer 诱导定理混同
%

\subsection{\texorpdfstring{交叉积与 \(2\)-余圈}{交叉积与 2-余圈}}
\label{b2:subsec:1804:01}

固定有限 Galois 扩张 \(K/k\)，记 \(G=\operatorname{Gal}(K/k)\)、\(n=|G|=[K:k]\)。希望为每个 \(g\in G\) 配一个基元 \(u_g\)，使它对标量的作用由 \(g\) 给出，而两个基元相乘只与 \(u_{gh}\) 相差一个非零标量，即
\[
u_g a=g(a)u_g,\qquad u_g u_h=c(g,h)u_{gh}.
\]
这两个规则把结合律转成一个可计算的条件：
\[
(u_g u_h)u_t=c(g,h)c(gh,t)u_{ght},\qquad
u_g(u_h u_t)=g\bigl(c(h,t)\bigr)c(g,ht)u_{ght}.
\]
若再要求 \(u_1=1\)，就得到下面的余圈方程及归一化条件。

\begin{definition}\label{b2:def:normalized-two-cocycle}
设 \(K/k\) 为有限 Galois 扩张，\(G=\operatorname{Gal}(K/k)\)，\(G\) 按域自同构作用于 \(K^\times\)。若函数 \(c:G\times G\to K^\times\) 对任意 \(g,h,t\in G\) 满足
\begin{equation}\label{b2:eq:two-cocycle}
c(g,h)c(gh,t)=g\bigl(c(h,t)\bigr)c(g,ht),\qquad
c(1,g)=c(g,1)=1,
\end{equation}
则称 \(c\) 为\term[guiyihua-er-yuquan]{归一化 \(2\)-余圈}[normalized 2-cocycle]。若取元素族 \((b_g)_{g\in G}\subseteq K^\times\) 且 \(b_1=1\)，定义函数
\[
(\delta b)(g,h)=b_g\cdot g(b_h)b_{gh}^{-1}
\]
并称 \(\delta b\) 为\term[er-yubian]{\(2\)-余边}[2-coboundary]。
\end{definition}

直接将 \(\delta b\) 代入\eqref{b2:eq:two-cocycle}，两边都化为 \(b_g\cdot g(b_h)gh(b_t)b_{ght}^{-1}\)，所以余边确为余圈。

\ProseParagraphBreak
余边记录的是同一组基元的缩放。若将 \(u_g\) 换成 \(u'_g=b_g u_g\)，其中 \(b_1=1\)，则
\[
u'_g u'_h
=b_g g(b_h)c(g,h)b_{gh}^{-1}u'_{gh}
=\bigl(c\delta b\bigr)(g,h)u'_{gh}.
\]
因此改变这些基元会改变乘法系数，但只使余圈乘上一个余边。

\begin{definition}\label{b2:def:galois-crossed-product}
设 \(K/k\) 为有限 Galois 扩张，\(G=\operatorname{Gal}(K/k)\)，\(c:G\times G\rightarrow K^\times\) 为归一化 \(2\)-余圈。在左 \(K\) 空间 \(B_c=\bigoplus_{g\in G}Ku_g\) 上规定
\begin{equation}\label{b2:eq:crossed-product-multiplication}
(a u_g)(b u_h)=a\cdot g(b)c(g,h)u_{gh}.
\end{equation}
按双可加性延拓后，所得含幺 \(k\) 代数称为由 \(K/k\) 和 \(c\) 确定的\term[jiaocha-ji-daishu]{交叉积代数}[crossed product algebra]。
\end{definition}
\symidx{\(B_c=(K/k,G,c)\)：Galois 扩张的交叉积代数}

其中 \(K\) 通过 \(a\mapsto au_1\) 嵌入；乘法公式同时规定了 \(u_g b=g(b)u_g\) 与 \(u_g u_h=c(g,h)u_{gh}\)。前一条控制基元经过标量时发生的自同构，后一条控制基元之间的乘法误差。前述换基计算也给出一个同构 \(B_{c\delta b}\to B_c\)：把新代数中指标 \(g\) 的基元送到 \(b_g u_g\) 即可。

\begin{proposition}\label{b2:prop:crossed-product-central-simple}
设 \(K/k\) 为有限 Galois 扩张，\(G=\operatorname{Gal}(K/k)\)，\(n=[K:k]\)，\(c:G\times G\rightarrow K^\times\) 为归一化 \(2\)-余圈，则相应交叉积 \(B_c=\bigoplus_{g\in G}Ku_g\) 是次数 \(n\) 的中心单 \(k\) 代数，单位为 \(u_1\)。子域 \(K\) 分裂它；每个 \(u_g\) 都可逆。
\end{proposition}
\begin{proof}
比较三项乘积 \((a u_g)(b u_h)(d u_t)\)：除共同因子 \(a\cdot g(b)gh(d)\) 外，两种结合方式的系数分别是\eqref{b2:eq:two-cocycle}的两边。因此乘法结合。归一化条件使 \(u_1\) 为单位，底域标量在中心。向量空间维数为 \(n[K:k]=n^2\)。由余圈条件还有 \(c(g,g^{-1})=g\bigl(c(g^{-1},g)\bigr)\)，故
\[
u_g^{-1}=c(g^{-1},g)^{-1}u_{g^{-1}}.
\]
左右相乘均为一。

中心的计算与循环代数同样，先检查与 \(K\) 交换的条件，再检查与全部基元交换的条件。若 \(x=\sum_g a_g u_g\) 与全部 \(b\in K\) 交换，则比较系数得 \(a_g\bigl(g(b)-b\bigr)=0\)。对 \(g\ne1\) 可选使差非零的 \(b\)，所以只有单位指标的系数可能非零。再与各 \(u_g\) 交换，就使该系数被全部 \(G\) 固定，因而属于 \(k\)。故中心是 \(k\)。

单性的证明沿用循环代数中减少非零项数的方法。对一个非零双边理想，选非零项数最少的元素，并取其中一个非零项的指标 \(g\)。右乘 \(u_g^{-1}\) 使每个指标 \(h\) 变为 \(hg^{-1}\)，各系数只乘上非零标量，所以非零项数不变，而单位指标的系数变为非零。再左乘一个 \(K^\times\) 中的标量，使该系数为一。所得元素与任意 \(b\in K\) 的交换子在单位指标处为零，其他非零项只可能来自原有的指标；若交换子非零，就有更少的非零项，违反最少性。于是该元素与 \(K\) 交换，由刚才的系数比较只能为一，理想即为全代数。这证明单性。最后 \([K:k]=n=\deg B_c\)，由\cref{b2:prop:maximal-field-splits-csa}，\(K\) 分裂 \(B_c\)。
\end{proof}

若交叉积分裂为 \(n\) 阶矩阵代数，它在一个 \(n\) 维 \(k\) 空间上的作用限制到 \(K\)，就成为一维 \(K\) 空间上的作用。在这个空间上，每个 \(u_g\) 的作用满足 \(T_g(az)=g(a)T_g(z)\)；选择一个 \(K\) 基以后，它只能是域自同构 \(g\) 再乘上一个非零标量。这些标量正是余边公式中的 \(b_g\)。

\begin{proposition}\label{b2:prop:crossed-product-split-coboundary}
设 \(K/k\) 为有限 Galois 扩张，\(G=\operatorname{Gal}(K/k)\)，\(c:G\times G\rightarrow K^\times\) 为归一化 \(2\)-余圈，\(B_c\) 为相应交叉积代数，则 \(B_c\) 分裂，当且仅当 \(c\) 是 \(2\)-余边。
\end{proposition}
\begin{proof}
若 \(c(g,h)=b_g\cdot g(b_h)b_{gh}^{-1}\)，其中 \(b_1=1\)，在 \(k\) 空间 \(K\) 上让 \(K\) 左乘，让 \(u_g\) 按 \(T_g(z)=b_g\cdot g(z)\) 作用。计算得到
\[
T_gT_h(z)=b_g g(b_h)gh(z)=c(g,h)T_{gh}(z),\qquad
T_g(az)=g(a)T_g(z).
\]
因此这些算子满足交叉积的两个乘法规则，且 \(T_1=\operatorname{id}_K\)，给出含幺同态 \(B_c\to\operatorname{End}_k(K)\)。由单性及两端同为 \(n^2\) 维，它是同构，其中 \(n=[K:k]\)。

若 \(B_c\cong\operatorname{End}_k(W)\)，由 \(\deg B_c=n\) 有 \(\dim_k W=n\)。限制到 \(K\) 后，维数公式给出 \(\dim_KW=1\)。选一个 \(K\) 基，将 \(W\) 识别为 \(K\)。由于 \(u_g z=g(z)u_g\)，其作用满足 \(T_g(z)=g(z)T_g(1)=b_g g(z)\)，其中 \(b_g=T_g(1)\ne0\)，且 \(b_1=1\)。代入 \(u_g u_h=c(g,h)u_{gh}\)，得 \(b_g\cdot g(b_h)=c(g,h)b_{gh}\)，正是余边条件。此时把 \(u_g\) 换成 \(b_g^{-1}u_g\)，它们在 \(K\) 上就按 \(g\) 作用，乘法误差也同时消失。
\end{proof}

\subsection{Brauer 群与群上同调}
\label{b2:subsec:1804:02}

\(2\)-余圈按逐点相乘组成交换群，余边构成其子群。定义商群
\[
H^2(G,K^\times)=\{\text{归一化 }2\text{-余圈}\}/\{2\text{-余边}\}.
\]
这个商群是当前 \(G\) 作用下二阶\term[qun-shangtongdiao]{群上同调}[group cohomology]的显式模型。第四卷第 12.1 节从 bar 消解建立余链复形，第 12.2 节再用低阶余圈解释群扩张；这里把交换群 \(K^\times\) 的运算写成乘法，得到同一套二阶方程，下面的论证只用这些显式方程。
\symidx{\(H^2(G,K^\times)\)：乘法系数的二阶群上同调}

\ProseParagraphBreak
对于 \(g\in G\)，可加映射 \(T_g\) 若满足 \(T_g(av)=g(a)T_g(v)\)，称其为 \(g\)-\term[banxianxing-yingshe]{半线性映射}[semilinear map]。在矩阵代数上，逐条目施加 \(g\) 就给出这种映射。下面把一个已分裂代数的下降数据写成具体矩阵。

\begin{proposition}\label{b2:prop:split-algebra-factor-set}
设 \(K/k\) 为有限 Galois 扩张，\(G=\operatorname{Gal}(K/k)\)，\(A\) 为由 \(K\) 分裂的中心单 \(k\) 代数，\(m=\deg A\)，取 \(K\) 代数同构 \(\rho:A_K=A\otimes_kK\rightarrow M_m(K)\)。令 \(\alpha_g=\rho(\operatorname{id}_A\otimes g)\rho^{-1}\) 为 \(A_K\) 的自然 Galois 作用在矩阵代数上的表示，\(g(X)\) 表示对矩阵 \(X\) 逐条目施加 \(g\)，则存在 \(C_g\in\operatorname{GL}_m(K)\)、\(C_1=I_m\)，使 \(\alpha_g(X)=C_g\cdot g(X)C_g^{-1}\)。由
\[
C_g\cdot g(C_h)=c(g,h)C_{gh}
\]
定义的标量 \(c(g,h)\in K^\times\) 构成归一化 \(2\)-余圈，而且 \([A]=-[B_c]\)，其中 \(B_c\) 是相应交叉积代数。若将 \(C_g\) 换为 \(b_gC_g\)，其中 \(b_g\in K^\times\)、\(b_1=1\)，则余圈换为 \(c\delta b\)。若将分裂模型换为 \(\rho'=\operatorname{Ad}(P)\rho\)，其中 \(P\in\operatorname{GL}_m(K)\)、\(\operatorname{Ad}(P)(X)=PXP^{-1}\)，并取
\[
C'_g=PC_g g(P)^{-1},
\]
则得到原余圈 \(c\)。特别地，所得类 \([c]\in H^2(G,K^\times)\) 与 \(C_g\) 及分裂同构 \(\rho\) 的选择无关。
\end{proposition}
\begin{proof}
\(\alpha_g\) 在标量上作用为 \(g\)，即为 \(g\)-半线性代数自同构。将它与逐条目作用 \(g^{-1}\) 依次复合，所得 \(\alpha_g\circ g^{-1}\) 是一个 \(K\) 代数自同构；由\cref{b2:thm:skolem-noether}，存在 \(C_g\in\operatorname{GL}_m(K)\)，使
\begin{equation}\label{b2:eq:brauer-semilinear-matrices}
\alpha_g(X)=C_g\cdot g(X)C_g^{-1},\qquad C_1=I_m.
\end{equation}
\(\alpha_g\alpha_h=\alpha_{gh}\) 说明 \(C_g\cdot g(C_h)C_{gh}^{-1}\) 与所有矩阵交换，所以为唯一的标量 \(c(g,h)\in K^\times\)。比较同一个矩阵积的两种结合方式，得到
\[
\begin{aligned}
\bigl(C_g g(C_h)\bigr)gh(C_t)
&=c(g,h)c(gh,t)C_{ght},\\
C_g g\bigl(C_h h(C_t)\bigr)
&=g\bigl(c(h,t)\bigr)c(g,ht)C_{ght}.
\end{aligned}
\]
消去可逆矩阵 \(C_{ght}\)，即为\eqref{b2:eq:two-cocycle}；\(C_1=I_m\) 又给出归一化条件。

对同一个 \(\alpha_g\)，两个实现矩阵只相差非零标量，因为它们的商与全部矩阵交换。因此任意另一组归一化的实现矩阵都形如 \(b_gC_g\)、\(b_1=1\)，而
\[
(b_gC_g)g(b_hC_h)
=b_g g(b_h)c(g,h)C_{gh}
=\bigl(c\delta b\bigr)(g,h)(b_{gh}C_{gh}).
\]
再改变分裂模型。由 \(\rho'=\operatorname{Ad}(P)\rho\) 得
\[
\alpha'_g(X)
=P C_g g(P)^{-1}g(X)g(P)C_g^{-1}P^{-1}.
\]
所述 \(C'_g\) 因而实现 \(\alpha'_g\)，且 \(C'_1=I_m\)。它们满足
\[
C'_g g(C'_h)
=PC_g g(C_h)gh(P)^{-1}
=c(g,h)C'_{gh},
\]
所以余圈不变。任意两个 \(K\) 代数分裂同构之差都是 \(M_m(K)\) 的 \(K\) 自同构，由 Skolem--Noether 为这种内自同构；这就涵盖了全部分裂模型的选择。

在 \(k\) 空间 \(V=K^m\) 上，让 \(K\) 按标量乘法作用，让 \(u_g\) 按 \(v\mapsto C_g\cdot g(v)\) 作用，得到 \(B_c\) 的表示。它含幺且 \(B_c\) 单，因而忠实。与 \(K\) 作用交换的算子恰好是 \(K\) 线性算子，可写为 \(X\in M_m(K)\)；这样的 \(X\) 与所有 \(u_g\) 作用交换，当且仅当
\[
XC_g=C_g g(X),\qquad\text{即 }X=\alpha_g(X)\quad(g\in G).
\]
这些不动矩阵恰为 \(\rho(A\otimes1)\)：在 \(A_K\) 中按 \(A\) 的一组 \(k\) 基展开，自然 Galois 作用的不动性恰好使全部系数属于 \(K^G=k\)。因此 \(B_c\) 在 \(\operatorname{End}_k(V)\) 中的中心化子同构于 \(A\)。由\cref{b2:cor:csa-centralizer-factorization}的中心为 \(k\) 的情形，
\[
B_c\otimes_k A\cong\operatorname{End}_k(V).
\]
所以 \([A]=-[B_c]\)，得到所述关系。
\end{proof}

\begin{theorem}\label{b2:thm:relative-brauer-cocycles}
设 \(K/k\) 为有限 Galois 扩张，\(G=\operatorname{Gal}(K/k)\)。记 \(H^2(G,K^\times)\) 为归一化 \(2\)-余圈按\(2\)-余边取的商群，\(\operatorname{Br}(K/k)=\ker\bigl(\operatorname{Br}(k)\rightarrow\operatorname{Br}(K)\bigr)\)，则交叉积给出群同构
\[
H^2(G,K^\times)\longrightarrow\operatorname{Br}(K/k),
\qquad[c]\longmapsto[B_c].
\]
\end{theorem}
\begin{proof}
首先，\cref{b2:prop:crossed-product-central-simple}使像落在相对 Brauer 群；换基公式保证它不依赖余圈代表。

为核对乘法，取两个余圈 \(c,d\)，并记 \(n=[K:k]\)。\(B_c\otimes_kB_d\) 的 \(k\) 维数为 \(n^4\)，而 \(B_{cd}\) 的维数为 \(n^2\)；当 \(n>1\) 时两者尺寸不同，一般不能期待它们直接同构。要比较的是 Brauer 类，需要去掉不影响除代数代表的矩阵大小；下面通过非零幂等元取角来实现这一点。有限 Galois 扩张具有代数同构
\begin{equation}\label{b2:eq:galois-tensor-diagonal}
K\otimes_kK\longrightarrow\prod_{g\in G}K,
\qquad a\otimes b\longmapsto\bigl(a\cdot g(b)\bigr)_{g\in G}.
\end{equation}
具体地，第一卷定理16.4.2给出本原元素 \(K=k(\theta)\)；其可分最小多项式在 \(K\) 中分成互异因子 \(T-g(\theta)\)。对 \(K[T]/(f)\) 应用第一卷定理10.4.3的中国剩余定理，便得到上述同构。

选择单位指标，是因为该坐标把 \(a\otimes b\) 送到 \(ab\)：两份 \(K\) 在这一分量上的作用被对齐。记 \(e_t\in K\otimes_kK\) 为在指标 \(t\) 处坐标为一、其他坐标为零的幂等元，并取 \(e=e_1\)，就有 \(e(a\otimes1)e=e(1\otimes a)e\)。因此希望取角以后只保留对两份 \(K\) 同步施加同一自同构的项。在 \(E=B_c\otimes_kB_d\) 中，用 \(u_g,v_h\) 表示两侧的基元，记 \(z_{g,h}=u_g\otimes v_h\)。共轭 \(z_{g,h}\) 对 \(K\otimes_kK\) 作用为 \(g\otimes h\)。若 \(\lambda_t(a\otimes b)=a\,t(b)\)，则
\[
\lambda_t\bigl((g\otimes h)(a\otimes b)\bigr)
=g(a)\,th(b)
=g\bigl(\lambda_{g^{-1}th}(a\otimes b)\bigr).
\]
按可加性，这个等式对 \(K\otimes_kK\) 的全部元素成立。将它用于 \(e\)，可知 \(z_{g,h}e z_{g,h}^{-1}=e_{gh^{-1}}\)。不同坐标幂等元的乘积为零，因而
\[
e z_{g,h}e=e e_{gh^{-1}}z_{g,h}=0\quad\text{若 }g\ne h.
\]
当 \(g=h\) 时，\(z_{g,g}\) 与 \(e\) 交换，令 \(w_g=e z_{g,g}e=e z_{g,g}\)。取角后，系数环 \(e(K\otimes_kK)e\) 经 \(\lambda_1\) 识别为 \(K\)。原代数是自由左 \(K\otimes_kK\) 模，以全部 \(z_{g,h}\) 为基；取角消去了 \(g\ne h\) 的项，留下
\[
eEe=\bigoplus_{g\in G}e(K\otimes_kK)w_g
\cong\bigoplus_{g\in G}K w_g
\]
的左 \(K\) 模分解。这里各项仍然独立，因为它们在原自由模中有不同的指标；而 \(w_g\ne0\)，否则右乘 \(z_{g,g}^{-1}\) 就会使 \(e=0\)。计算乘法时，\(g\otimes g\) 在单位坐标上正是 \(g\)，并且
\[
w_g w_h=e\bigl(c(g,h)\otimes d(g,h)\bigr)z_{gh,gh}
=c(g,h)d(g,h)w_{gh}.
\]
同样，若用 \(e(b\otimes1)e\) 表示角代数中的标量 \(b\in K\)，便有
\[
w_g b=e(g(b)\otimes1)z_{g,g}=g(b)w_g.
\]
于是把 \(B_{cd}\) 的基元送到 \(w_g\)，把标量送到 \(e(K\otimes_kK)e\) 中的对应元素，得到含幺 \(k\) 代数同构 \(B_{cd}\to eEe\)，角代数的单位是 \(e=w_1\)。由\cref{b2:prop:brauer-division-representative}，中心单代数的非零幂等角与原代数有同一 Brauer 类。因此 \([B_c]+[B_d]=[B_{cd}]\)，映射保持群运算。\(c=1\) 时的交叉积由\cref{b2:prop:crossed-product-split-coboundary}分裂，确实对应零元。

若 \([B_c]=[B_d]\)，由\cref{b2:prop:brauer-division-representative}，可写 \(B_c\cong M_r(D)\)、\(B_d\cong M_s(D)\)，其中 \(D\) 为同一个中心除代数。两代数的次数同为 \(n\)，所以 \(r\deg D=s\deg D=n\)，从而 \(r=s\)，可取 \(k\) 代数同构 \(\phi:B_c\to B_d\)。记 \(\iota_c,\iota_d\) 为两代数中 \(K\) 的标准嵌入。由\cref{b2:thm:skolem-noether}，有 \(z\in B_d^\times\)，使
\[
\iota_d(a)=z\phi\bigl(\iota_c(a)\bigr)z^{-1}\qquad(a\in K).
\]
于是 \(\psi=\operatorname{Ad}(z)\phi\) 在两份标准 \(K\) 之间为恒等。将 \(x=\psi(u_g)\) 展开为 \(\sum_h b_{g,h}v_h\)，由 \(xa=g(a)x\) 得
\[
b_{g,h}h(a)=g(a)b_{g,h}\qquad(a\in K).
\]
若 \(b_{g,h}\ne0\)，可消去这个 \(K\) 中的系数，得到 \(h(a)=g(a)\) 对全部 \(a\in K\) 成立，故 \(h=g\)。因此 \(\psi(u_g)=b_gv_g\)，且因 \(u_g\) 可逆，有 \(b_g\in K^\times\)；单位条件给出 \(b_1=1\)。比较 \(\psi(u_g u_h)\) 与 \(\psi(u_g)\psi(u_h)\) 的系数，得到
\[
c(g,h)b_{gh}=b_g\cdot g(b_h)d(g,h).
\]
因此 \(c/d\) 为余边，证明单射性。

最后，对任意 \([A]\in\operatorname{Br}(K/k)\)，\(A_K\) 的 Brauer 类为零，由\cref{b2:prop:brauer-division-representative}在底域 \(K\) 上的结论，它同构于一个 \(K\) 上的全矩阵代数。因此可应用\cref{b2:prop:split-algebra-factor-set}，得到余圈 \(c\)，使 \([A]=-[B_c]\)。由已经证明的乘法相容性，\([B_{c^{-1}}]=-[B_c]\)，所以原类是 \([c^{-1}]\) 的像。这证明满射性，所需群同构成立。
\end{proof}

半线性映射作为 \(k\) 线性算子参与交叉积的表示，却同时改变 \(K\) 标量。乘法公式中的 \(g\bigl(c(h,t)\bigr)\) 正记录了这种改变。从余圈构造交叉积，与从分裂代数的半线性作用提取因子集，这两种构造的关系为
\[
[c]\longmapsto[B_c],\qquad
A\longmapsto c_A,\qquad
[A]=-[B_{c_A}]=[B_{c_A^{-1}}].
\]
前一个箭头是\cref{b2:thm:relative-brauer-cocycles}的同构；从一个已分裂的代数 \(A\) 出发按\cref{b2:prop:split-algebra-factor-set}构造因子集时，原类对应的是 \([c_A^{-1}]\)。

\begin{corollary}[循环情形的范数商]\label{b2:cor:cyclic-relative-brauer}
设 \(K/k\) 为次数 \(n\ge2\) 的循环 Galois 扩张，\(G=\operatorname{Gal}(K/k)=\langle\sigma\rangle\)，则有群同构
\[
H^2(G,K^\times)\cong\operatorname{Br}(K/k)
\cong k^\times/\operatorname{N}_{K/k}(K^\times).
\]
具体地，映射 \(a\operatorname{N}_{K/k}(K^\times)\mapsto[(K/k,\sigma,a)]\) 给出范数商到相对 Brauer 群的同构。
\end{corollary}
\begin{proof}
由\cref{b2:thm:relative-brauer-cocycles}，每个相对 Brauer 类都由某个交叉积 \(B_c\) 表示。在其中取 \(u=u_\sigma\)。反复使用乘法规则，可知 \(u^i\) 是 \(u_{\sigma^i}\) 的非零 \(K\) 倍数，特别地 \(u^n=a\in K^\times\)。因为 \(u\) 与其幂 \(a\) 交换，而 \(ua=\sigma(a)u\)，可消去 \(u\) 得 \(\sigma(a)=a\)，所以 \(a\in k^\times\)。将原基换成 \(1,u,\ldots,u^{n-1}\)，就是一次前述缩放基元的操作；它不改变余圈的上同调类，却将乘法误差化为
\[
u^i u^j=
\begin{cases}
u^{i+j},&i+j<n,\\
a\,u^{i+j-n},&i+j\ge n,
\end{cases}
\qquad(0\le i,j<n).
\]
因此相应余圈为
\[
c_a(\sigma^i,\sigma^j)=a^{\lfloor(i+j)/n\rfloor},
\]
且 \(B_c\cong(K/k,\sigma,a)\)。这证明参数给出从 \(k^\times\) 到 \(\operatorname{Br}(K/k)\) 的满射。进位余圈满足 \(c_a c_b=c_{ab}\)，由\cref{b2:thm:relative-brauer-cocycles}，该参数映射是群同态。

若再用 \(u'=bu\)、\(b\in K^\times\)，则
\[
(bu)^n=b\,\sigma(b)\cdots\sigma^{n-1}(b)\,u^n
=\operatorname{N}_{K/k}(b)a.
\]
因而乘以一个范数正对应再次改变这组基元。由\cref{b2:thm:cyclic-algebra-norm-splitting}，参数映射的核恰为 \(\operatorname{N}_{K/k}(K^\times)\)。商群同构定理遂给出所述范数商，连同\cref{b2:thm:relative-brauer-cocycles}便得到全部同构。
\end{proof}

交叉积的基元除了记录乘法，也使约化迹可以直接计算。左乘 \(u_h\) 将指标 \(g\) 移到 \(hg\)，因此只有单位指标的项会贡献矩阵的对角元；计算时仍须把左 \(K\) 系数转换成右 \(K\) 坐标。

\begin{proposition}\label{b2:prop:crossed-product-reduced-trace}
设 \(K/k\) 为有限 Galois 扩张，\(G=\operatorname{Gal}(K/k)\)，\(c:G\times G\to K^\times\) 为归一化 \(2\)-余圈，\(B_c=\bigoplus_{g\in G}Ku_g\) 为相应交叉积代数。对任意 \(x=\sum_{g\in G}a_g u_g\)、\(a_g\in K\)，有
\[
\operatorname{Trd}_{B_c/k}(x)=\operatorname{Tr}_{K/k}(a_1)
=\sum_{g\in G}g(a_1).
\]
\end{proposition}
\begin{proof}
记 \(n=[K:k]\)，将 \(B_c\) 的右 \(K\) 空间记为 \(W\)。由 \(a u_g=u_g g^{-1}(a)\)，\((u_g)_{g\in G}\) 也是 \(W\) 的一组基。根据\cref{b2:prop:maximal-field-splits-csa}中的分裂同构，
\[
B_c\otimes_k K\longrightarrow
\operatorname{End}_K(W),\qquad
x\otimes b\longmapsto(y\mapsto xyb),
\]
其中目标作用在原代数的右 \(K\) 空间上，而来源是标量扩张代数。由于 \(\dim_KW=n\)，右端是 \(M_n(K)\) 的一个模型。在这个模型中，\(x\otimes1\) 是左乘 \(x\) 的算子。对 \(h,g\in G\)，
\[
(a_hu_h)u_g
=a_h c(h,g)u_{hg}
=u_{hg}(hg)^{-1}\bigl(a_h c(h,g)\bigr).
\]
因此在列 \(g\) 中，这一项的矩阵系数只出现在行 \(hg\)。若 \(h\ne1\)，行指标不等于列指标，故它不贡献对角元；若 \(h=1\)，由 \(c(1,g)=1\)，对角元恰为 \(g^{-1}(a_1)\)。左乘 \(x\) 的矩阵迹于是为
\[
\sum_{g\in G}g^{-1}(a_1)=\sum_{g\in G}g(a_1)
=\operatorname{Tr}_{K/k}(a_1).
\]
有限 Galois 扩张的域迹等于全部共轭之和，而\cref{b2:def:reduced-trace-norm}中的约化迹就是任意分裂模型里的矩阵迹，故得所述公式。
\end{proof}

第五卷第 5.3 节的 Brauer 诱导定理研究有限群特征标，与本章按张量积运算的 Brauer 群是不同对象。
